% =============================================================
%  Section: Experimentation
% =============================================================

This section traces six successive training runs (V1--V6) of the YOLOv1
implementation on COCO~2017. The first four runs all produced
$\mathrm{mAP}@0.5 = 0$ on the held-out test set despite training losses that
appeared to converge. The recovery from this state of affairs is the central
narrative of this work: a sequence of hyper-parameter changes that failed to
diagnose the problem, followed by a code audit that identified two
implementation bugs invisible to the unit-test suite. V5 and V6 were trained
after the fixes and produced the first non-zero detection scores.

We adopt the convention that each version is identified by its short tag
(V1\ldots V6) and we cross-reference the corresponding row of
Table~\ref{tab:run-summary} throughout.

% -------------------------------------------------------------
\subsection{Common training setup}
\label{sec:experiments:setup}

All six versions share the optimiser, the loss formulation (whose
implementation, however, differed in V1--V4 from V5--V6 due to the bug fixes
discussed in Section~\ref{sec:bug-investigation}), the data pipeline, and the
evaluation protocol. We summarise the constants here and only describe
per-version variations in their respective subsections.

\begin{itemize}
    \item \textbf{Dataset:} COCO~2017~\citep{lin2014microsoft}, all 80 classes,
        consolidated into a single annotation file and split deterministically
        (seed $42$) into 90\,\%\,/\,5\,\%\,/\,5\,\% for train, validation and
        test. The test split contains 6\,165 images and 42\,795 ground-truth
        bounding boxes after the removal of crowd annotations. Of these, the
        dataset class retains only $31\,167$ at training time because the
        one-object-per-cell rule (Section~\ref{sec:impl:dataset}) drops any
        annotation whose grid cell is already claimed by a previously
        encountered object; the dropped boxes still count as missed
        detections at evaluation, since the evaluator reads the raw COCO
        annotations.
    \item \textbf{Optimiser:} stochastic gradient descent with Nesterov-free
        momentum~$0.9$ and weight decay~$5\!\times\!10^{-4}$, applied uniformly
        to all parameters (no parameter groups).
    \item \textbf{Learning-rate schedule:} a linear warm-up of three epochs
        from $0$ to the peak rate, followed by cosine annealing
        \citep{loshchilov2017sgdr} down to a minimum of
        $1\!\times\!10^{-5}$. The peak rate is the per-version
        parameter and is reported individually.
    \item \textbf{Mixed precision:} the forward pass runs under
        \texttt{torch.amp.autocast} in float16, while master weights and the
        loss accumulator remain in float32~\citep{micikevicius2018mixed}.
        Gradient scaling is performed with \texttt{GradScaler} and the scaled
        gradients are unscaled before clipping.
    \item \textbf{Distributed setup:} the training loop is written for
        \texttt{DistributedDataParallel} with \texttt{torchrun}, but in
        practice every run used a single GPU and a per-GPU batch size of
        $32$. V1--V3 ran on a RunPod pod equipped with $8\!\times$ NVIDIA H100
        GPUs (of which only one was effectively exercised due to the storage
        bottleneck described in Section~\ref{sec:impl:training});
        V4--V6 ran on a Vast.ai H200 NVL pod ($140$~GB HBM3) with
        node-local NVMe.
    \item \textbf{Loss:} the five-term sum-of-squared-errors loss from
        \citet{redmon2016you}, with $\lambda_{\text{coord}}=5$ and
        $\lambda_{\text{noobj}}=0.5$. The implementation is the object of
        Section~\ref{sec:bug-investigation}.
    \item \textbf{Evaluation:} the trained model is run on the full test split
        with prediction probability thresholds of $0.1$ and $0.3$, followed by
        per-class non-maximum suppression at IoU~$0.5$. The reported metric is
        mean Average Precision at IoU~$0.5$, computed per class using the
        unclaimed-GT bookkeeping described in
        Section~\ref{sec:theory:map} and averaged across the 80 classes that have at
        least one ground-truth instance in the test split (which in our case is
        all of them).
\end{itemize}

The variables that change across versions are the backbone, the peak learning
rate, the gradient clipping threshold, the number of training epochs and, most
importantly for V5 and V6, the loss and decoding code. They are summarised in
Table~\ref{tab:run-summary}.

\begin{table}[ht]
\centering
\caption{Summary of all six training runs. ``Bugs'' indicates whether the loss
and coordinate-decoding bugs analysed in Section~\ref{sec:bug-investigation}
were present at training time. Validation loss is the final per-batch
sum-reduced value from the training log; the magnitude jump between V4 and V5
is itself evidence of the bug fix and is discussed in
Section~\ref{sec:experiments:v5}.}
\label{tab:run-summary}
\small
\begin{tabular}{lllcrrcc}
\toprule
\textbf{Run} & \textbf{Backbone} & \textbf{Bugs} & \textbf{Epochs} &
\textbf{Train loss} & \textbf{Val loss} &
\textbf{mAP@0.5 (0.1)} & \textbf{mAP@0.5 (0.3)} \\
\midrule
V1 & Darknet24 & present & 33 (diverged) & --     & N/A  & 0.0000 & 0.0000 \\
V2 & Darknet24 & present & 100           & 1.45   & 1.50 & 0.0000 & 0.0000 \\
V3 & Darknet24 & present & 135           & 1.20   & 1.35 & 0.0000 & 0.0000 \\
V4 & ResNet18  & present & 50            & 0.96   & 0.98 & 0.0000 & 0.0000 \\
V5 & ResNet18  & fixed   & 50            & 336.34 & 337.37 & 0.0409 & 0.0373 \\
V6 & ResNet18  & fixed   & 60            & 330.50 & 337.40 & \textbf{0.0415} & 0.0381 \\
\bottomrule
\end{tabular}
\end{table}

% -------------------------------------------------------------
\subsection{V1: paper-faithful Darknet baseline}
\label{sec:experiments:v1}

The first attempt follows \citet{redmon2016you} as literally as the constraints
of a modern PyTorch~\citep{paszke2019pytorch} stack allow. The backbone is a
re-implementation of the paper's 24-layer Darknet, initialised from scratch
with the default PyTorch initialisation. The detection head is two fully
connected layers (50\,176~$\to$~4096~$\to$~4410) with dropout~0.5 after the
first \citep{srivastava2014dropout}. The peak learning rate is
$1\!\times\!10^{-3}$, matching the paper's choice for training from scratch on
ImageNet-free data. No gradient clipping is applied. The intended schedule is
100 epochs.

Training diverges at epoch~33: the loss tensor turns into NaN within a single
batch and remains NaN for every subsequent step. The proximate cause is the
combination of a randomly initialised 24-layer convolutional network, a
high peak learning rate, and the lack of any clipping mechanism: a single
unlucky batch produces a gradient norm large enough to push the master weights
into a region where the next forward pass overflows the fp16 dynamic range.
The run is abandoned, and no checkpoint is selected for evaluation. We
nonetheless include V1 in Table~\ref{tab:run-summary} for narrative
completeness.

% -------------------------------------------------------------
\subsection{V2: regularisation by gradient clipping}
\label{sec:experiments:v2}

V2 is a direct response to V1: the peak learning rate is halved to
$5\!\times\!10^{-4}$ and we introduce gradient-norm clipping with a threshold
of $5.0$. The architecture is unchanged. The model trains for the planned
100 epochs without divergence; the validation loss bottoms out around 1.50 and
plateaus. On the test set, however, the mAP@0.5 is still exactly zero. The
qualitative behaviour is that of a model that has learnt to drive object
confidences to almost zero everywhere: post-NMS detections are scarce and
their bounding boxes are unrelated to the ground-truth objects.

We interpreted this at the time as a sign that the lowered learning rate was
preventing the network from escaping a trivial minimum of the no-object
confidence term. In hindsight, this interpretation was incomplete: the
``trivial minimum'' was indeed real, but its existence was caused by the loss
bug discussed in Section~\ref{sec:bug:loss-indices}, not by the
hyper-parameters.

% -------------------------------------------------------------
\subsection{V3: relaxing the clip, restoring the LR}
\label{sec:experiments:v3}

V3 returns the peak learning rate to $1\!\times\!10^{-3}$ and increases the
gradient-clip threshold to $10.0$, with the working hypothesis that V2 was
under-trained. The number of epochs is extended to 135 to match the paper's
total schedule length. The training loss continues to decrease and reaches its
lowest value across the Darknet runs (1.20), but the validation loss begins to
diverge from training loss around epoch~$\sim 40$, which we interpreted as
overfitting. mAP@0.5 on the test set is again exactly zero.

At this point the working hypothesis became that the Darknet backbone, trained
from scratch on COCO without the ImageNet pre-training stage that
\citet{redmon2016you} performed, simply lacks the inductive bias needed for
this dataset. We therefore moved on to V4.

% -------------------------------------------------------------
\subsection{V4: ResNet18 pretrained on ImageNet}
\label{sec:experiments:v4}

V4 replaces the 24-layer Darknet with a ResNet18~\citep{he2016deep} loaded with
ImageNet weights from \texttt{torchvision} (\texttt{ResNet18\_Weights.DEFAULT}).
The classification head of the ResNet is discarded and the convolutional trunk
up to and including \texttt{layer4} produces feature maps of shape
$(B, 512, 14, 14)$ given the standard $448\times 448$ input. We bridge to the
detection head's expected $(B, 1024, 7, 7)$ with a two-layer adapter (a
stride-2 convolution from $512$ to $1024$ channels, followed by a stride-1
convolution that refines the representation), each with batch
normalisation~\citep{ioffe2015batch} and a LeakyReLU activation. The detection
head is identical to V1--V3. All parameters of the model carry
\texttt{requires\_grad=True} from epoch one; we do not freeze any layer.

Because the backbone now expects ImageNet-normalised inputs, the dataset
transform is augmented with the standard channel mean/std normalisation. The
peak learning rate is reduced to $1\!\times\!10^{-4}$ to fine-tune the
pretrained convolutional features without destroying them, while still leaving
enough learning capacity for the randomly initialised adapter and FC head
(which between them carry the majority of the model's 248~M parameters).
Gradient clipping is kept at $10.0$, and the schedule is shortened to 50
epochs.

V4 trains without numerical issues. The validation loss reaches 0.98 at the
best epoch---the lowest value yet, and on its face suggestive of a model that
has fit the data well. Yet the mAP@0.5 on the test set is again exactly zero.
Post-NMS we observe a different qualitative failure mode from V2/V3: the
model emits a large number of low-confidence boxes per image (on the order
of 20 with $p > 0.005$), but they are spatially uncorrelated with the ground
truth, and their predicted classes do not match.

This was the result that forced the question: \emph{is the loss decreasing
because the model is learning, or because it is finding a trivial minimum of a
loss that is not what we think it is}? The investigation that followed is the
subject of the next section.

% =============================================================
\subsection{Bug investigation}
\label{sec:bug-investigation}

The trigger for an audit was the observation that, across four runs with very
different hyper-parameter settings and two different backbones, the test mAP
remained identically zero while the validation losses spanned the range
$0.96$--$1.50$. A loss that decreases consistently across hyper-parameter
sweeps but produces a detector that detects nothing is suspect: either the
model genuinely has no localisation capacity in this regime (which is hard to
reconcile with a fully fine-tuned ResNet18 at val\_loss~$<\!1$), or the loss is
not enforcing the constraint we expect.

The audit identified two independent bugs, each of which is sufficient on its
own to produce $\mathrm{mAP}@0.5 = 0$. The first is in the loss function; the
second is in the post-processing routine that converts model outputs back to
image-relative bounding boxes. We discuss them in turn.

% -------------------------------------------------------------
\subsubsection{Bug~1: VOC-layout indices applied to COCO tensors}
\label{sec:bug:loss-indices}

The original implementation of the YOLO loss in \texttt{src/loss.py} addressed
the per-cell prediction tensor by hardcoded integer indices that matched the
Pascal VOC layout (20 classes, $C=20$). The expected per-cell layout for
$C=20$ is
\[
  \underbrace{[\,0,\;\ldots,\;19\,]}_{\text{class probabilities}}\;\;
  \underbrace{[\,20\,]}_{\text{conf}_1}\;\;
  \underbrace{[\,21,\;\ldots,\;24\,]}_{x_1,\,y_1,\,w_1,\,h_1}\;\;
  \underbrace{[\,25\,]}_{\text{conf}_2}\;\;
  \underbrace{[\,26,\;\ldots,\;29\,]}_{x_2,\,y_2,\,w_2,\,h_2},
\]
which the loss reads literally as \texttt{[..., 20:21]} for the
``object exists'' mask, \texttt{[..., 21:25]} for the first predicted box, and
so on. For COCO with $C=80$, the actual per-cell tensor produced by the
dataset is
\[
  \underbrace{[\,0,\;\ldots,\;79\,]}_{\text{class probabilities}}\;\;
  \underbrace{[\,80\,]}_{\text{conf}_1}\;\;
  \underbrace{[\,81,\;\ldots,\;84\,]}_{x_1,\,y_1,\,w_1,\,h_1}\;\;
  \underbrace{[\,85\,]}_{\text{conf}_2}\;\;
  \underbrace{[\,86,\;\ldots,\;89\,]}_{x_2,\,y_2,\,w_2,\,h_2},
\]
so the indices used by the loss point inside the \emph{class probability}
region of the tensor, never reaching the coordinate slots.

The downstream consequences are catastrophic, but each consequence is locally
silent:
\begin{enumerate}
    \item \texttt{target[..., 20:21]}, which the loss interprets as the
        object-presence indicator (\texttt{exists\_box}), is in COCO the
        class-probability slot for class index~20 (which after the dense COCO
        remapping is the \emph{elephant} class). The mask is therefore $1$
        only in cells whose ground-truth class is elephant, and $0$ otherwise.
    \item The coordinate slots that the loss reads
        (\texttt{predictions[..., 21:25]}, \texttt{target[..., 21:25]}) are
        actually class probability slots for class indices~21--24
        (\emph{bear}, \emph{zebra}, \emph{giraffe}, \emph{backpack}). The IoU
        computation between two slices of class probabilities returns a
        nonsense quantity that is essentially uncorrelated with anything the
        network can learn from.
    \item The class loss is computed over the slice
        \texttt{predictions[..., :20]}, so only the first twenty COCO classes
        (out of eighty) participate in the class term at all. Classes 20--79
        receive no gradient signal whatsoever.
    \item The no-object loss penalises \texttt{predictions[..., 20]} and
        \texttt{predictions[..., 25]} in cells where the target's position~20
        is zero. Position~20 of the target is zero in the overwhelming
        majority of cells (those that do not happen to contain an elephant),
        so the model is strongly pulled towards driving these two output
        positions to zero everywhere. Position~20 and~25 in the prediction
        tensor, however, are not the confidence slots; they are the
        class-probability slots for the class \emph{elephant} and (in
        practice, after \texttt{argmax}) the slot interpreted as ``conf~2''
        is the class probability for \emph{train}. The trivial-minimum
        behaviour observed in V2--V4 is a faithful learning of this confused
        no-object signal, not a property of the MSE loss as such.
\end{enumerate}

The bug is invisible from a unit-testing perspective. The existing test in
\texttt{tests/test\_loss.py} instantiates the loss with $C=20$:

\begin{lstlisting}[caption={Pre-fix unit test that passed for the wrong reason.},label=lst:voc-test]
loss_fn = YoloLoss(S=7, B=2, C=20)
target = torch.zeros(1, 7, 7, 30)
target[0, 3, 4, 0]    = 1                              # class
target[0, 3, 4, 20]   = 1                              # confidence
target[0, 3, 4, 21:25] = torch.tensor([0.5, 0.5, 0.3, 0.3])
\end{lstlisting}

\noindent At $C=20$, the indices $20$, $21\!:\!25$, $25$, $26\!:\!30$ that the
loss reads happen to coincide with the slots that the dataset writes, so the
test passes. The test would have to be parametrised on $C$ to detect the
mismatch; ours was not.

\paragraph{Verification.} We constructed a synthetic COCO-shaped target with
one ground-truth object at class~7, cell~$(3,4)$, with box
$(x_c, y_c, w, h) = (0.5, 0.5, 0.3, 0.3)$ and computed the loss on two
predictions:

\begin{enumerate}
    \item A ``perfect'' prediction that matches the target at all slots,
        including the correct COCO coordinate slots $[81\!:\!85]$.
    \item The same prediction but with the four box-coordinate slots
        $[81\!:\!85]$ overwritten with arbitrarily large values
        ($99.0,\ldots,99.0$).
\end{enumerate}

Under a correctly indexed loss, prediction~(2) must produce a much larger
value than prediction~(1). Under the buggy loss, the two values were
identical to within numerical precision: $0.000490$ in both cases. This
confirms that the loss was \emph{gradient-blind} to the slots that the model
is supposed to predict.

Aggregated over a sample of 200 real COCO training images, the buggy
$\texttt{exists\_box}$ mask was $1$ in only 3 cells out of $1\,062$ cells that
actually contained a ground-truth object---about $0.3\,\%$---and these three
were all elephants. The remaining $99.7\,\%$ of real objects contributed
exactly zero to the box and object terms of the loss.

\paragraph{Fix.} We replaced every hardcoded literal in the loss with an index
derived from \texttt{self.C}:

\begin{lstlisting}[caption={Index variables introduced at the top of
    \texttt{YoloLoss.forward}.}]
C          = self.C
OBJ_IDX    = C
BOX1_SLICE = slice(C + 1, C + 5)
CONF2_IDX  = C + 5
BOX2_SLICE = slice(C + 6, C + 10)
\end{lstlisting}

\noindent and propagated these names through every site that previously held a
literal. The class term was also corrected from
\texttt{predictions[..., :20]} to \texttt{predictions[..., :C]} so that all 80
classes participate. The pre-existing $C=20$ unit test continues to pass
because the new variables evaluate to the same numeric positions when $C=20$.
A full side-by-side diff is provided in Appendix~\ref{app:loss-diff}.

A second, smaller follow-up to the loss was committed shortly afterwards: the
square-root transform of $w$ and $h$ (the paper's trick to equalise the
contribution of small and large boxes) was implemented as an \emph{in-place}
assignment, \texttt{box\_predictions[..., 2:4] = \ldots}. In-place writes on a
tensor that participates in the autograd graph silently sever the backward
pass through that tensor, with no exception raised at runtime. The fix is to
re-assemble the modified tensor with \texttt{torch.cat} along the channel
dimension. We mention this here for completeness; it did not affect the V1--V4
runs because, given Bug~1, no gradient flowed through that path anyway.

% -------------------------------------------------------------
\subsubsection{Bug~2: spurious \texorpdfstring{$1/S$}{1/S} factor in \texorpdfstring{$w,\,h$}{w, h} decoding}
\label{sec:bug:wh-decoding}

The post-processing routine \texttt{convert\_cellboxes} in
\texttt{src/utils.py} converts the raw per-cell prediction tensor into a list
of bounding boxes in image-relative coordinates. The conversion of the
$(x, y)$ centre coordinates is straightforward: the predictions are stored as
offsets within the cell, so the global $x$ coordinate is the cell column index
plus the in-cell offset, divided by $S$:
\[
  x_{\text{global}} = \frac{j + x_{\text{cell}}}{S}, \qquad
  y_{\text{global}} = \frac{i + y_{\text{cell}}}{S}.
\]
The width and height, however, are stored by the dataset as ratios of the
\emph{whole image} dimensions ($w_{\text{norm}} = w_{\text{px}}/W_{\text{img}}$
and likewise for $h$), and the loss compares the model's $w,h$ output
\emph{directly} against these values. The model therefore learns to predict
$w_{\text{norm}}, h_{\text{norm}}$.

In our implementation the decoding code, however, performed an additional
division by $S$:
\begin{lstlisting}
w_h = (1.0 / S) * best_boxes[..., 2:4]   # BUG: divides w,h by 7
\end{lstlisting}

\noindent This is the symmetric mistake of porting from a different
encoding. The implementation we drew partial inspiration
from~\citep{persson2020yolo} stores width and height as $w_{\text{norm}}\cdot S$
and $h_{\text{norm}}\cdot S$ (``number of cells the box spans''), so dividing
by $S$ at decode time is correct \emph{in that codebase}. In ours it is not.

\paragraph{Consequence.} Every predicted box is decoded with width and height
seven times smaller than the value the model has learnt to output. Given that
the network was trained to predict $w_{\text{norm}}, h_{\text{norm}} \in [0,1]$,
this scaling produces boxes of at most one-seventh of the correct linear
dimension and roughly one-forty-ninth of the correct area.

The arithmetic upper bound on the IoU between such a shrunken prediction and
its perfectly aligned ground truth is
\[
  \text{IoU} \le \frac{(\tfrac{1}{S})^2}{1 + (\tfrac{1}{S})^2 - (\tfrac{1}{S})^2}
              = \frac{1}{S^2} \approx 0.020,
\]
which is far below the threshold of $0.5$ used in mAP@0.5. The bug is
therefore sufficient on its own to drive the metric to exactly zero, even in
the hypothetical scenario where the loss bug had not also existed and the
network had learnt to localise objects correctly.

\paragraph{Verification.} We encoded a known box with $w=h=0.3$ centred at
cell~$(3,3)$, passed the tensor through \texttt{convert\_cellboxes}, and
inspected the decoded width and height. The buggy code returned
$w_{\text{dec}} = h_{\text{dec}} = 0.0429$, which is exactly $0.3/7$, while
$x$ and $y$ round-tripped correctly. After removing the $1/S$ factor on
$w_h$, the round-trip is exact to floating-point precision.

\paragraph{Fix.} The fix is one line: drop the $1/S$ factor on $w_h$ while
keeping it on $x$ and $y$. Variable renamed from \texttt{w\_y} to \texttt{w\_h}
for clarity.

% -------------------------------------------------------------
\subsubsection{Why the unit-test suite did not catch either bug}
\label{sec:bug:tests}

Both bugs are detectable with a single round-trip check, but the existing test
suite covered shapes and value ranges only. The pre-existing tests for
\texttt{convert\_cellboxes} were:
\begin{itemize}
    \item \texttt{test\_convert\_cellboxes\_output\_shape}: asserts the output
        is of shape $(B, S^2, 6)$.
    \item \texttt{test\_convert\_cellboxes\_class\_indices\_in\_range}: asserts
        that the predicted class index column lies in $[0, C-1]$.
\end{itemize}

\noindent Neither inspects the value of $w, h$. A test of the form ``encode a
known box, decode it, assert that the recovered $w, h$ match'' would have
detected Bug~2 trivially. Likewise, a loss test parameterised on $C$ with
$C \in \{20, 80\}$ would have detected Bug~1 on the very first run: the
``perfect prediction'' test at $C=80$ would have returned a non-zero loss for
any prediction whose values at the bogus VOC slots did not happen to coincide
with the class probabilities at those slots. We added both round-trip tests
together with the fixes, as documented in
Section~\ref{sec:experiments:v5}.

\paragraph{Generalisable lesson.} The deeper failure here is not the specific
indices but the assumption that a unit test that compiles, runs and passes
is evidence that the function it tests is correct. The two tests above did
exactly what they claimed and provided zero evidence about the function's
mathematical correctness. In a detection pipeline, where intermediate tensors
have many slots and most of them are unused by any single test image, the
density of meaningful coverage per unit test is low. The round-trip checks
we added pin down a single quantitative invariant per function and break
when the function's mathematical contract changes, which is the level of
specificity that is missing from shape-only assertions.

% -------------------------------------------------------------
\subsection{V5: applying the fixes}
\label{sec:experiments:v5}

V5 retains every hyper-parameter from V4---the same ResNet18 backbone, the
same adapter, the same FC head, the same peak learning rate of
$1\!\times\!10^{-4}$, the same gradient-clip threshold of $10.0$, the same
50-epoch schedule, the same optimiser, batch size and warm-up
configuration---and changes only the loss and decoding code as described in
Section~\ref{sec:bug-investigation}. The point of holding everything else
equal is to attribute any change in the result to the fixes alone.

The first observation, available from the very first epoch of training, is
that the magnitude of the loss is now in a completely different regime.
Whereas V4 produced training and validation losses around $0.96$--$0.98$, V5
produces values in the range of several hundred. The final training and
validation losses are $336.34$ and $337.37$ respectively. We had expected this:
the loss is reduced with \texttt{nn.MSELoss(reduction="sum")} over all 90
slots of every cell of every image in a batch. In V1--V4, due to Bug~1, almost
all the per-cell summands collapsed to zero, because they were multiplied by
the ``exists\_box'' mask which was itself zero almost everywhere. The
non-trivially zero contributions came almost entirely from a few elephant
cells. Once the indices are corrected, every cell with an object contributes
to the box and confidence terms, and every cell contributes to the no-object
term and the class term. The absolute magnitude of the loss therefore rises
by two orders of magnitude. This is not a regression; it is the loss finally
seeing the data.

At evaluation time, V5 produces $\mathrm{mAP}@0.5 = 0.0409$ at
prob\_threshold~$0.1$ on the test set, and $0.0373$ at prob\_threshold~$0.3$
(see Table~\ref{tab:run-summary}). This is the first non-zero detection score
of the entire project. It is a low absolute number---we discuss in detail in
Section~\ref{sec:results} why this is so, and in the Discussion section what
would be required to push it higher---but it is a categorical change from the
preceding four runs.

The qualitative behaviour also changes. Detection boxes now cluster on actual
objects in the scene, predicted classes correlate with ground-truth classes on
the easy classes (animals with distinctive silhouettes, vehicles), and per-NMS
detection counts drop to a sensible 2--6 boxes per image at
prob\_threshold~$0.3$. We provide examples drawn from the test set, sorted by
qualitative quality, in Section~\ref{sec:results}.

% -------------------------------------------------------------
\subsection{V6: continuation of V5}
\label{sec:experiments:v6}

V6 is a continuation of V5 rather than a separate experiment. We resumed
training from the V5 ``best'' checkpoint and trained for an additional 10
epochs, holding every other setting identical. The motivation was to verify
that the V5 result was not a coincidence of the cosine schedule landing on a
favourable epoch, and to give the model a bit more time near the bottom of the
learning rate ramp.

The final validation loss is essentially identical to V5 (337.40 vs.\ 337.37),
indicating that the schedule has reached saturation. The test mAP improves
marginally to $0.0415$ at prob\_threshold~$0.1$ and $0.0381$ at $0.3$. We treat
V6 as the headline run of the project and use its checkpoint
(\texttt{v6\_run/best\_v6.pth}) for all subsequent analyses in
Section~\ref{sec:results}, including the full per-class AP table in
Appendix~\ref{app:per-class-ap} and the qualitative figures.

% -------------------------------------------------------------
\subsection{Loss curves and a note on the magnitude shift}
\label{sec:experiments:curves}

Figure~\ref{fig:loss-curves} plots the training and validation loss curves
across V4, V5 and V6 on a shared time axis. The visual jump in absolute
magnitude between V4 and V5 is the empirical signature of the index fix. The
right panel zooms in on V5+V6 to make the slow descent of both training and
validation losses visible at the appropriate scale.

\begin{figure}[!ht]
\centering
\includegraphics[width=0.95\textwidth]{loss_comparison_v3.png}
\caption{Training and validation loss across runs. Left: V4 (bugs present)
plateauing near $1.0$. Middle: V5+V6 (fixes applied) plateauing near $337$.
Right: zoom on the last ten epochs of V5+V6, showing the slow improvement that
the absolute scale makes invisible in the middle panel. The magnitude jump
between V4 and V5 is itself the signature of the bug fix discussed in
Section~\ref{sec:bug-investigation}.}
\label{fig:loss-curves}
\end{figure}

A reader who reaches V5 in the training log without the context of
Section~\ref{sec:bug-investigation} could reasonably mistake the rise in loss
for a regression. We emphasise that the comparison between V4 and V5 is a
comparison between two losses that compute fundamentally different
quantities, not between a well-trained and a poorly-trained version of the
same loss.
